\documentclass{article}
\usepackage{geometry}
\usepackage{amsmath}
\usepackage{minibox}
\title{Pre-Algebra to Algebra 1 Summer Work}
\author{OLMC Math Department}
\date{\today}
\begin{document}
\maketitle
\newpage\
\begin{center}
	\title{Mental Operations}
\end{center}
\paragraph{Complete the following operations as quickly as possible. \textbf{\underline{Do Not Use A Calculator.}} Focus on speed as well as accuracy}
\section{Addition of Whole Numbers}

\textbf{Students whould be able to quickly add two digit whole numbers in their heads. Try to keep the time for each problem under 5 seconds.}
\bigskip


\framebox{\begin{minipage}{0.85 \textwidth}

    \textbf{Tip:} While traditional Addition algorithms usually involve moving from right to left and carrying ones. 
    It is often faster to mentally move from left to right by turning one difficult problem into three easy problems.

    \bigskip
    \textbf{Example:} To solve $97+68$ mentally, I first calculate $90+60=150$ and set the result aside. I then calculate $7+6=13$ and combine the two results: $13+150=163$

\end{minipage}
}
\begin{align*}
        27&+16 \qquad 14+39 \qquad 36+42 \qquad 91+54 \qquad 64+38 \qquad 43+14 \qquad 57+66 \\\\
        32&+59 \qquad 79+4  \qquad 76+88 \qquad 43+18 \qquad 38+82 \qquad 63+23 \qquad 57+18 \\\\
        23&+64 \qquad 81+21 \qquad 99+48 \qquad 37+28 \qquad 24+24 \qquad 78+8  \qquad 33+55 \\\\
        82&+56 \qquad 94+57 \qquad 70+34 \qquad 26+12 \qquad 39+78 \qquad 45+32 \qquad 12+86 \\\\
        87&+58 \qquad 80+90 \qquad 72+21 \qquad 63+21 \qquad 72+60 \qquad 83+26 \qquad 13+75 \\\\
        49&+63 \qquad 30+28 \qquad 99+30 \qquad 37+62 \qquad 22+64 \qquad 29+12  \qquad 63+77 \\\\
        49&+11 \qquad 90+73 \qquad 79+30 \qquad 8+47  \qquad 69+37 \qquad 57+21 \qquad 12+65 \\\\
        22&+79 \qquad 44+89 \qquad 21+85 \qquad 65+87 \qquad 48+36 \qquad 14+18 \qquad 47+12 \\\\
\end{align*}

\newpage


\begin{center}
	\title{Mental Operations (cont.)}
\end{center}
\paragraph{Complete the following operations as quickly as possible. \textbf{\underline{Do Not Use A Calculator.}} Focus on speed as well as accuracy}

\section{Addition of terminating decimals}

\textbf{The inclusion of terminating decimals should not significantly impact the time to solve addition problems}
\bigskip


\framebox{\begin{minipage}{0.85 \textwidth}

    \textbf{Tip:} The strategy for adding decimals that most students are familiar with is to "line up" the decimals. Another option is to multiply both addends by a power of ten that removes the decimals, perform the addition and than divide the result by the same power of ten. It is also a good idea to always double check that your answer is "reasonable." Your result should never be more than 2 away from the sum of the addend's whole number parts.

    \bigskip
	\textbf{Example:} To solve $3.69+2.4$ mentally, I first multiply both addends by 100 and calculate $369+240=609$. Then I simply divide the result by 100: $609/100=6.09$. Since 6.09 is between 3 and 7 [$(3+2)\pm2$], the result is considered "reasonable."

\end{minipage}}

\begin{align*}
	4.37&+6.6 \qquad  0.001+0.02 \qquad 0.003+0.08  \qquad  1.2+0.1  \qquad 0.006+0.5  \qquad  5.5+3.8   \\\\
	0.01&+0.04  \qquad  5.1+2.72   \qquad   7.4+4.34  \qquad  0.3+0.8  \qquad   2.2+1.4  \qquad  5.7+0.46 \\\\
	2.8&+2.57   \qquad  7.3+0.8    \qquad   5.2+9.68  \qquad 0.07+0.06 \qquad  0.05+0.03 \qquad 3.52+0.4  \\\\
	1.4&+4.7    \qquad  9.2+9.5    \qquad 0.008+0.006 \qquad  3.2+2.3  \qquad  0.09+0.08 \qquad  7.7+7.2   \\\\
	6.5&+2.5    \qquad  2.6+3.5    \qquad  0.05+0.09  \qquad 0.04+0.05 \qquad   6.8+1.8  \qquad  6.1+6.5   \\\\
	4.6&+1.777  \qquad 1.66+5.3    \qquad    7.9+1.07    \qquad   4.4+8.2   \qquad    0.007+0.9   \qquad  9.8+7.09  \\\\
\end{align*}

\newpage

\begin{center}
	\title{Mental Operations (cont.)}
\end{center}
\section{Subtraction of Whole Numbers}

\textbf{Students whould be able to quickly subtract two digit whole numbers in their heads. Try to keep the time for each problem under 5 seconds.}
\bigskip


\framebox{\begin{minipage}{0.85 \textwidth}

    \textbf{Tip:} Just like with addition we should move right to leftand convert one big problem into three small ones.

    \bigskip
	
	\textbf{Example:} To solve $63-38$ mentally, I first round the subtrahend (subtracted number) up and perform the subtraction operation: $63-40=23$. Then I take my rounded subtrahend and subtract the original number: $40+38=2$.  Finally, I combine the two results: $23+2=25$
\end{minipage}}

\begin{align*}
        80&-58 \qquad 97-19 \qquad 84-42 \qquad 91-17 \qquad 64-41 \qquad 78-24 \qquad 97-49 \\\\
        54&-28 \qquad 94-50 \qquad 84-15 \qquad 82-56 \qquad 98-39 \qquad 48-29 \qquad 84-22 \\\\
        90&-13 \qquad 73-17 \qquad 74-33 \qquad 71-48 \qquad 85-8  \qquad 57-9  \qquad 82-37 \\\\
        97&-72 \qquad 91-50 \qquad 69-45 \qquad 70-32 \qquad 91-36 \qquad 72-24 \qquad 77-45 \\\\
        80&-14 \qquad 66-35 \qquad 89-10 \qquad 95-61 \qquad 54-20 \qquad 82-15 \qquad 92-43 \\\\
        40&-11 \qquad 98-70 \qquad 89-43 \qquad 67-33 \qquad 95-30 \qquad 49-24 \qquad 76-14 \\\\
        68&-12 \qquad 89-37 \qquad 85-37 \qquad 99-12 \qquad 61-5  \qquad 91-47 \qquad 35-14 \\\\
        76&-39 \qquad 73-22 \qquad 32-7  \qquad 87-46 \qquad 73-29 \qquad 96-19 \qquad 85-27 \\\\
\end{align*}

\newpage
\begin{center}
	\title{Mental Operations (cont.)}
\end{center}

\paragraph{Complete the following operations as quickly as possible. \textbf{\underline{Do Not Use A Calculator.}} Focus on speed as well as accuracy}

\section{Subtraction of terminating decimals}

\textbf{The inclusion of terminating decimals should not significantly impact the time to solve subtraction problems}
\bigskip


\framebox{\begin{minipage}{0.85 \textwidth}

    \textbf{Tip:} You can use the same multiplying by a power of 10 and then dividing strategy we used for addition ofterminating decimals.

	\bigskip
	
	\textbf{Example:} To solve $1.3-0.6$ mentally, I first multiply both addends by 10 and calculate $13-6=7$. Then I simply divide the result by 10: $7/10=0.7$. Since 0.7 is between -1 and 3 [$(1+0)\pm2$], the result is considered "reasonable."

\end{minipage}}

\begin{align*}
	7.7&-4.0 \qquad  9.6-1.4 \qquad 5.9-2.3  \qquad  0.92-0.56  \qquad 0.53-0.05  \qquad  0.89-0.41  \\\\
	0.39&-0.17  \qquad  1-0.29   \qquad   9.2-4.1  \qquad  4.9-1.6  \qquad   10-6.4  \qquad  0.49-0.17 \\\\
	0.92&-0.64   \qquad  0.99-0.46    \qquad   0.96-0.14  \qquad 0.78-0.24 \qquad  0.93-0.29 \qquad 0.58-0.11  \\\\
	0.7&-0.25    \qquad  0.5-0.04    \qquad 8-5.6 \qquad  3.2-1.1  \qquad  4.1-1.3 \qquad  0.77-0.23  \\\\
	4.6&-1.5    \qquad  7.6-2.2    \qquad  4.9-2.3  \qquad 0.84-0.48 \qquad   9.2-0.3  \qquad  0.71-0.35 \\\\
	9.0&-6.2   \qquad 0.98-0.35    \qquad    0.81-0.13   \qquad   6.6-3.7   \qquad    8.4-5.7   \qquad  0.88-0.13 \\\\
\end{align*}

\newpage

\end{document}
